\section{本章小结}
\label{sec:dep-summary}

本章从$X\sim\operatorname{Unif}(-1,1)$与$Y=X^{2}$出发：协方差为零，依赖却是确定的。协方差只检查一对线性函数，独立要求联合分布等于边缘乘积，即对所有有界函数对乘积公式都成立。依赖度量就是对这一差别的量化，评价它时要分开三件事：零值是否刻画独立，非零数值随依赖怎样变化，以及能否从有限样本可靠判断。本章只讨论前两件。

互信息是联合分布相对边缘乘积的KL散度，零值在任何可测空间上刻画独立。它在双向可测的一一变换下不变，在不可逆的局部处理下不增，连续边缘时只依赖Copula。它不看取值之间的距离，也没有有限上界：在标准Borel空间上，只要$Y=f(X)$且$Y$的分布无原子，联合分布就集中在边缘乘积给零概率的图像上，互信息为无穷大。有限的计算下界不等于真实值。

距离协方差用权重$\lVert\bm{t}\rVert_2^{-(p+1)}\lVert\bm{s}\rVert_2^{-(q+1)}$对特征函数之差作平方积分。一阶矩保证它有限，零值由特征函数的唯一性刻画独立。二阶矩有限时它等于成对距离的三项期望，也等于双重中心化距离乘积的期望；中心化消去了只由边缘决定的部分，而指数$1$（不是$2$）使它看到非线性依赖。用距离方差归一化得到取值于$[0,1]$的距离相关，它在平移、正交变换和整体缩放下不变，取$1$当且仅当两组坐标只差一个相似变换，常量情形约定为$0$。

HSIC是联合分布与边缘乘积在乘积核下的平方MMD，等于核期望的三项组合，也等于交叉协方差算子的Hilbert--Schmidt范数平方。乘积核为特征核时零值刻画独立，高斯核满足这一条件；线性核只剩协方差的平方，距离诱导核给出距离协方差的四分之一。高斯核的带宽改变数值，甚至改变两个分布的次序。

离散化评分把取值装进有限个格子。固定划分的标签互信息不超过原互信息，可能恰好错过依赖；沿加细的划分它逼近互信息。在允许任意切点的网格中搜索并按$\ln\min(k,l)$归一化，总体评分的零值刻画实值变量的独立，分段单调的无噪声函数关系得分为$1$。经验网格评分是样本的函数，小样本下对独立变量也会给出明显的正值。

这些量能同时判别独立，却按不同的方式度量依赖：互信息量信息的减少，距离相关量距离的同步，HSIC量指定尺度上的特征协方差，网格评分量网格能否捕捉关系。选择之前要说清楚随机对象所在的空间、可积条件、需要保持的变换，并把零值判别与非零强度比较分开。用样本估计这些量、校准估计误差和检验独立性，由卷八的非参数推断与假设检验两章承接；算出一个样本分数，并不等于完成了统计判断。
